\subsection*{用运动规律识别一个函数}
从复平面的点一出发，让速度永远等于当前位置乘以 $\mathrm i$。速度与半径垂直，长度与半径相同，所以它沿单位圆以单位角速度运动。走过时间 $t$ 后的位置应是 $\cos t+\mathrm i\sin t$。

另一边，指数函数 $\mathrm e^{\mathrm it}$ 对时间的导数是 $\mathrm i\mathrm e^{\mathrm it}$，起点也是一。两个表达式满足同一微分方程和初值，在相应唯一性结论下，它们就是同一条轨迹。欧拉公式因此可被理解为两种描述相遇。

\subsection*{特殊角度是整个运动的路标}
时间 $\pi/2$ 到达 $\mathrm i$，时间 $\pi$ 到达负一，时间 $2\pi$ 回到一。著名恒等式是半圈位置的记录；若只背五个常数排在一起，就会丢失背后的完整连续运动。

在动画中暂停于四分之一圈和半圈，分别读横纵分量，再把它们与余弦正弦对应。改变转动方向只改变到达路径，走正半圈或负半圈都到负一，这也展示了指数的周期性。

周期性将在复对数中变成多值性：同一个终点可以对应许多角度，所以从终点反推指数不能保持全局唯一。欧拉公式既统一了指数与三角，也提醒我们把实数反函数的习惯推广到复数时需要重新检查。
